\section{Pinned cycles and the return walk}\label{sec:loops}

We first express a pinned spin deficit as a return probability for one
particle.  The resulting walk moves through a periodic set of holes.
We state precisely the loop identities and pin estimates imported from
\cite{Ordering,BlochLaw}; the estimates that improve their precision are
proved in Section~\ref{sec:geometry}.

\subsection{The picture and its pin law}

Fix a finite free-boundary cube $\Lambda=\Lambda_N\subset\Z^3$ and write
$H_{S,\Lambda}=H_{S,N}$ and $\ell=2S$.  Its slot set is
$V_\Lambda=\Lambda\times\{1,\ldots,\ell\}$.
A slot edge joins every pair of slots whose sites are nearest neighbors.
The base picture consists of independent Poisson clocks of rate $1/2$
on these undirected edges during $[0,\beta)$, and independent uniform
permutations of the $\ell$ slots at each site at the seam $0$.
Following the swaps and then the seam permutations produces directed
cycles on the time circle $\R/\beta\Z$.  A slot-time point at an
instantaneous operation is understood to lie just after that operation.
Write $\EE_{\mathrm b}$ for expectation under the base law.
This is the exchange-cycle representation of T\'oth \cite{Toth1993},
with the higher-spin symmetrization discussed by Nachtergaele
\cite[Section~2.2]{Nachtergaele1994}.

Let $\partial_{\mathrm{in}}\Lambda$ be the sites of $\Lambda$ having a
neighbor outside it, and put
\[
 D_0=\partial_{\mathrm{in}}\Lambda\times\{1,\ldots,\ell\}
                              \times(\R/\beta\Z).
\]
An \emph{allowed pin set} is $D=D_0\cup F$, where $F$ is a finite
deterministic set of slot-time points.  Pins require the corresponding
cycles to be up.  If $n_D(\omega)$ is the number of cycles missing $D$,
define
\begin{equation}\label{eq:pinlaw}
 Z_D=\EE_{\mathrm b}2^{n_D},\qquad
 \PP_D(d\omega)=Z_D^{-1}2^{n_D(\omega)}\PP_{\mathrm b}(d\omega),
 \qquad
 \mathcal B_D=\bigcup_{C\cap D=\varnothing} C.
\end{equation}
Under the associated color law, cycles meeting $D$ are up and all other
cycles are independently up or down with equal probabilities.  These
definitions also make sense for $D=\varnothing$.  We denote probabilities
under the color law by $\PP_D^{\mathrm{col}}$.

Hereafter all low-temperature bounds are uniform in the finite cube and
in its allowed deterministic pin set.  Their constants may depend on
the fixed spin and on fixed geometric constants, but not on the number
or placement of the extra pins.

\subsection{Exposure and a fresh walk}

For an allowed $E$ and an unpinned point $i$, set $P=E\cup\{i\}$.
Sample the picture under $\PP_P$ and expose all cycles in $\mathcal B_P$,
including their inter-site jumps and their seam mappings.  At each time
their occupied slots are called \emph{obstacles}; the complementary slots
are \emph{holes}.  Every pin is a hole at its prescribed time.

At fixed exposure, the fresh walk on holes has rate $1/2$ on every edge
between holes.  When an obstacle swaps into a hole, that hole is
transported in the opposite direction.  A swap of two obstacles has no
effect on the holes.  At the seam the walk uses a fresh uniform bijection
between the remaining slots at each site, respecting the exposed part of
the seam permutation.  The exposure is repeated periodically, while all
transition randomness of the walk is resampled on each turn.  This is
a time-inhomogeneous Markov walk at fixed exposure; it is not obtained
by following the unexposed part of a single conditioned picture.

Let $q_i$ be its probability, starting from $i$, of returning to the
same slot-time point at strictly positive elapsed time before hitting
$E$.  A return can occur only at an integer number of periods.  We use
$r=\EE_Pq_i$ and $d_E(i)=\PP_E^{\mathrm{col}}(i\text{ is down})$.

\begin{proposition}[Loop and pin inputs]\label{prop:loopinputs}
The following statements hold at every finite volume and $\beta>0$.
\begin{enumerate}
\item If $q_\Lambda$ is the number of slot edges and $c(\omega)$ the
number of cycles, then
\[
 \Tr e^{-\beta H_{S,\Lambda}}
       =e^{\beta q_\Lambda/4}\EE_{\mathrm b}2^{c(\omega)}.
\]
The spectral distribution of the total $z$-spin is obtained by the
unpinned color law, assigning $+1/2$ to each up slot and $-1/2$ to each
down slot at time zero.  The expectation, in this normalized color law,
of the exponential of $\beta$ times the field times this sum is the
ratio of the field and zero-field partition functions.
\item For the allowed sets $E,P$ and point $i$ above,
\begin{equation}\label{eq:pinreturn}
 \PP_E(i\in\mathcal B_E)=\frac{2r}{1+r},
 \qquad d_E(i)=\frac{r}{1+r},\qquad r=\EE_Pq_i.
\end{equation}
\item More generally, let $T$ be a nonempty finite set of points at one
deterministic time, disjoint from an allowed $D$, and put $P'=D\cup T$.
Expose $\mathcal B_{P'}$ and let $q_j^{D,T}$ be the fresh-walk probability
from $j\in T$ of hitting $T$ at positive elapsed time before $D$.  Then
\begin{equation}\label{eq:simultaneouspin}
 \PP_D(T\subset\mathcal B_D)\le (2r_T)^{|T|},
 \qquad
 r_T=\frac1{|T|}\sum_{j\in T}\EE_{P'}q_j^{D,T}.
\end{equation}
\item For any allowed $D$, the factorial measure of $j$ distinct
inter-site marks under $\PP_D$ is bounded by $2^j$ times the base
factorial measure.  This holds with chronological ordering and with
arbitrary deterministic restrictions on the edges and times.
\end{enumerate}
\end{proposition}

The first assertion is the exact spin-$S$ representation in
\cite[Proposition~2.2]{Ordering}, with the present rate and Hamiltonian
normalizations.  The second combines the one-pin weight identity and
fresh-turn identity in \cite[Propositions~2.4 and~2.8]{BlochLaw}.
The third combines \cite[Proposition~3.1 and Lemma~4.4]{Ordering}; in
particular, no spatial separation of the test points is assumed.
The last assertion is \cite[Lemma~2.4]{Ordering}.  The statements include
all-time boundary pins by the dense-boundary limit proved there.

\subsection{Kernel conventions and energy}

Kernels act forward on column vectors, with counting inner products and
norms.  Killing on a pin means a projection at its prescribed time;
boundary killing is imposed throughout, including at interval endpoints.
Killing on exit from a spatial set retains the exit rates on the
generator's diagonal.  The adjoint of an interval reverses the order of
all factors, inverts transports, transposes the seam averages, and
reverses the order of killing projections.  Thus the two sides of an
instantaneous operation remain distinct slices.  Each seam factor is
included exactly once in a traversal.

Extend the hole graph by making every slot outside $\Lambda$ available.
Its edges join available slots over neighboring sites.  For a vector
extended by zero on the inner boundary and outside $\Lambda$, define
\begin{equation}\label{eq:dirichlet}
 \D_t(v)=\frac12\sum_{\{a,b\}\text{ available at }t}
                                      (v(a)-v(b))^2.
\end{equation}
The exterior extension will allow flows to infinity; all test vectors
will have finite support.

\begin{lemma}[Contraction and energy loss]\label{lem:kernelenergy}
At fixed exposure, every killed interval kernel and its adjoint are
substochastic in both counting sums and contract counting $\ell^p$ for
$p=1,2,\infty$.  Between instantaneous operations a propagated vector
satisfies
\[
 \frac{d}{dt}\norm{v_t}_2^2=-2\D_t(v_t).
\]
The squared norm cannot increase at an instantaneous operation.

Let $U$ be a one-period kernel with any allowed pin killings, and let $y$
vanish on its cut killing slots.  Set
$e_y=\ip{y}{(I-U)y}\ge0$.  For either $U_d=U$ or $U^*$, let $Y_t$ be
the adjoint propagation of $y$ from the terminal cut to an intermediate
slice.  Then
\begin{equation}\label{eq:energy}
 \int_{\mathrm{traversal}}\D_t(Y_t)\,dt\le e_y,
 \qquad
 \norm{y-U_d^*y}_2^2\le2e_y.
\end{equation}
The contraction and local norm-loss statements hold with extra spatial
killing and for intervals continued across a period cut.
Equation~\eqref{eq:energy} concerns one complete traversal; each copy
of that traversal has the same bound.
\end{lemma}

\begin{proof}
For one realization of the fresh exchanges and seam completions, the
unkilled trajectories define a bijection between the endpoint hole sets.
Deleting trajectories that encounter a killing set gives a partial
permutation.  Averaging proves the two substochastic bounds and the
$\ell^1$ and $\ell^\infty$ contractions.  Rowwise Cauchy--Schwarz gives
the $\ell^2$ contraction, as in
\cite[Section~3.1]{BlochLaw}.  Transposition gives the stated reverse
realization and the same bounds.

On a smooth interval the generator is symmetric, with rate $1/2$ on
each available edge and Dirichlet killing rates.  Summing by edges proves
the norm-loss identity.  Zero extension gives exactly
\eqref{eq:dirichlet}: a surviving interior slot has no edge directly to
the exterior of the cube.  Instantaneous contractions can only add to
the norm loss.  Finally, for $V=U$ or $U^*$, the real quadratic forms
agree and
\[
 \norm{y}_2^2-\norm{Vy}_2^2
       =2e_y-\norm{y-Vy}_2^2.
\]
This identity, contraction, and the integrated norm loss prove
\eqref{eq:energy}.
\end{proof}

All these are finite-volume statements with piecewise smooth evolution.
We disregard the null event of coincident Poisson marks or marks at a
prescribed deterministic time.  Deterministic operations that coincide
are composed in their assigned order.  The kernels are measurable from
finite-time transition formulas, and hitting probabilities from
increasing finite-horizon limits.  Consequently we may average them
under the actual picture law, including on events defined below.
